\documentclass[12pt]{article}
\usepackage[T1]{fontenc}
\usepackage[utf8]{inputenc}
\usepackage{lmodern}
\usepackage{textcomp}
\usepackage{enumitem}
\usepackage{geometry}
\usepackage{graphicx}
\usepackage{amsmath, amssymb}
\geometry{margin=1in}
\begin{document}
\begin{center}
Economics - Undergraduate Level Assessment 4 \
Total Marks: 40 \
Answer all questions.
\end{center}

\section*{Multiple Choice (10 marks)}
\begin{enumerate}[label=\arabic*.]
\item If the CPI goes to 150 from 120 then prices have
\begin{enumerate}[label=(\alph*)]
    \item fallen 50 percent.
    \item fallen 25 percent.
    \item risen 15 percent.
    \item risen 20 percent.
    \item fallen 20 percent.
\end{enumerate}
\item A small business estimates price elasticity of demand for the product to be 3. To raise total revenue, owners should
\begin{enumerate}[label=(\alph*)]
    \item keep the price constant as demand is elastic.
    \item decrease price as demand is inelastic.
    \item decrease the quality of the product.
    \item increase price as demand is elastic.
    \item increase price as demand is inelastic.
\end{enumerate}
\item Economics supposedly deals with the efficient allocation of scarce goods. Where does scarcity fit in the circular flow model?
\begin{enumerate}[label=(\alph*)]
    \item Scarcity is relevant only for non-renewable resources in the circular flow model
    \item Scarcity is only relevant in the factor markets of the circular flow model
    \item Scarcity is only considered in the household sector of the circular flow model
    \item Scarcity is addressed through government intervention in the circular flow model
    \item Scarcity is implicit in the entire analysis of the circular flow model
\end{enumerate}
\item According to the quantity theory of money increasing the money supply serves to
\begin{enumerate}[label=(\alph*)]
    \item increase long-run output with no significant impact on inflation.
    \item have no significant impact on either output or inflation.
    \item lower the unemployment rate while also lowering the rate of inflation.
    \item increase short-run output but it is the source of long-run inflation.
    \item increase the nation's long-run capacity to produce.
\end{enumerate}
\item How might the traditional view of economic efficiency in an oligopolyand the historical facts surrounding many oligopol-isticmanufacturing industries disagree with each other?
\begin{enumerate}[label=(\alph*)]
    \item Historical data contradicts theory
    \item Oligopolies are less common in practice than in theory
    \item There is no contradiction between theory and practice
    \item Historical data is insufficient to determine the efficiency of oligopolies
    \item Oligopolies typically demonstrate higher innovation rates than predicted
\end{enumerate}
\end{enumerate}

\section*{True / False (10 marks)}
\textit{Answer True or False}
\begin{enumerate}[label=\arabic*.]
\setcounter{enumi}{5}
\item True or False: If nominal GDP is $6000 and the GDP deflator is 200, then real GDP equals $6000.
\item True or False: Raising taxes on disposable income is an effective measure to combat built-in inflation by reducing consumer spending.
\item True or False: Social overhead capital primarily consists of private sector investments in technology startups aimed at innovation and profit.
\item True or False: Lower interest rates can lead to decreased private investment, which enhances the impact of expansionary fiscal policy.
\item True or False: In the long run, a tax cut in a full employment economy will have no effect on the price level but will increase real output.
\end{enumerate}

\section*{Long Form Response (20 marks)}
\textit{Answer each question with a clear and complete explanation.}
\begin{enumerate}[label=\arabic*.]
\setcounter{enumi}{10}
\item Discuss the relationship between income elasticity of demand for real money and real income growth. Given an elasticity of 0.8 and a demand growth of 2.6\%, calculate and analyze the growth rate of real income.
\item Discuss how public debt can create inflationary pressures in an economy, focusing on its impact on consumption spending and the implications of delayed purchasing power.
\end{enumerate}

\vfill
\noindent\textit{End of Paper}
\end{document}